% =====================================================================
\chapter{Evaluación, calibración y seguridad}
\label{tema:6}
% =====================================================================

Este tema estudia la conducta de un modelo desde fuera: cómo medir su rendimiento con conjuntos de pruebas, cómo cuantificar la incertidumbre de esas medidas, cuándo su confianza es fiable, por qué alucina y cómo se buscan y se miden sus fallos de seguridad. La organización sigue la clase~11, «Benchmarking and Evaluation», de CS224N~\cite{cs224n} y la semana de robustez adversaria de CS~2881r de Harvard~\cite{cs2881}. Los conjuntos de evaluación proceden de Hendrycks \emph{et al.}~\cite{hendrycks}, Liang \emph{et al.}~\cite{helm} y Chen \emph{et al.}~\cite{chen}; la calibración, de Guo \emph{et al.}~\cite{guo} y Kadavath \emph{et al.}~\cite{kadavath}; las alucinaciones, de Kalai \emph{et al.}~\cite{kalai25}, y los ataques, de Zou \emph{et al.}~\cite{zou} y Ganguli \emph{et al.}~\cite{ganguli}. Los intervalos de confianza siguen la clase~23 del curso 18.05 del MIT~\cite{mit1805}.

\section{Conjuntos de evaluación y su incertidumbre}

Un conjunto de evaluación (\emph{benchmark}) es una colección fija de tareas con un criterio automático de acierto.
\begin{itemize}[leftmargin=1.4em]
  \item \emph{MMLU}~\cite{hendrycks} reúne preguntas de opción múltiple de 57 materias, desde matemáticas elementales hasta derecho; en su publicación, el mayor GPT-3 superaba el azar en casi 20 puntos porcentuales de media.
  \item \emph{HELM}~\cite{helm} evalúa 30 modelos en 16 escenarios básicos con siete métricas (exactitud, calibración, robustez, equidad, sesgo, toxicidad y eficiencia), además de 26 escenarios específicos, y eleva la fracción de combinaciones de modelo y escenario evaluadas del $17{,}9\,\%$ al $96{,}0\,\%$.
  \item \emph{HumanEval}~\cite{chen} contiene 164 problemas de programación escritos a mano, cada uno con su firma, su descripción y pruebas unitarias; una solución es correcta si supera todas las pruebas.
\end{itemize}

\begin{definicion}[Exactitud empírica]\label{def:exactitud}
Sea un modelo con probabilidad de acierto $\theta$ en una población de tareas, y sean $n$ tareas extraídas de forma independiente. Sea $X_k\in\{0,1\}$ el acierto en la tarea $k$. La \emph{exactitud empírica} es $\bar X=\frac1n\sum_kX_k$.
\end{definicion}

\begin{proposicion}[Intervalo conservador]\label{prop:ic}
Con la notación de la \cref{def:exactitud}, $\E[\bar X]=\theta$ y $\Var(\bar X)=\theta(1-\theta)/n\le 1/(4n)$. En consecuencia, por la aproximación normal, $\bar X\pm z_{\alpha/2}/(2\sqrt n)$ es un intervalo de confianza aproximado de nivel al menos $1-\alpha$ para $\theta$; con $z_{0{,}025}\approx 2$, el intervalo al $95\,\%$ es $\bar X\pm 1/\sqrt n$.
\end{proposicion}

\begin{proof}
Las $X_k$ son variables de Bernoulli independientes de parámetro $\theta$, con media $\theta$ y varianza $\theta(1-\theta)$; por la linealidad de la esperanza y la independencia, $\E[\bar X]=\theta$ y $\Var(\bar X)=n\,\theta(1-\theta)/n^{2}$. La función $\theta(1-\theta)$ es una parábola con máximo $1/4$ en $\theta=1/2$. Por el teorema central del límite (módulo~5), $(\bar X-\theta)/\sqrt{\theta(1-\theta)/n}$ es aproximadamente normal estándar, y sustituir la desviación por su cota superior $1/(2\sqrt n)$ solo ensancha el intervalo~\cite[clase~23, §3]{mit1805}.
\end{proof}

Con $n=1000$ preguntas, el margen al $95\,\%$ es $1{,}96/(2\sqrt{1000})=0{,}031$: dos modelos con exactitudes del $70\,\%$ y del $72\,\%$ no se distinguen con esa muestra (\cref{eje:6-mit}). Las diferencias de uno o dos puntos que se publican a menudo exigen conjuntos de varios miles de preguntas, o comparaciones emparejadas sobre las mismas preguntas. A la incertidumbre estadística se suma la \emph{contaminación} (módulo~22): si las preguntas aparecieron en los datos de entrenamiento, la exactitud mide memoria y no capacidad.

\section{pass@\texorpdfstring{$k$}{k} y su estimador insesgado}

\begin{definicion}[pass@$k$]\label{def:passk}
Sea un problema que el modelo resuelve en cada intento, de forma independiente, con probabilidad $p$. Su \emph{pass@$k$} es la probabilidad de que al menos una de $k$ muestras sea correcta, $1-(1-p)^{k}$. El pass@$k$ de un conjunto de problemas es la media de los pass@$k$ de sus problemas~\cite[§2.1]{chen}.
\end{definicion}

Calcular $1-(1-\hat p)^{k}$ con una estimación $\hat p$ de $p$ da un valor sesgado (\cref{eje:6-sesgo}). Chen \emph{et al.} generan $n\ge k$ muestras por problema, cuentan las $c$ correctas y usan otro estimador.

\begin{proposicion}[Estimador insesgado de pass@$k$]\label{prop:passk}
Sean $n\ge k$ muestras independientes, cada una correcta con probabilidad $p$, y sea $c$ el número de correctas. Entonces
\[
\widehat{\mathrm{pass}}_k=1-\frac{\binom{n-c}{k}}{\binom nk}
\]
cumple $\E\bigl[\widehat{\mathrm{pass}}_k\bigr]=1-(1-p)^{k}$, con el convenio $\binom{n-c}{k}=0$ si $n-c<k$. Además, si $n-c\ge k$,
\[
\frac{\binom{n-c}{k}}{\binom nk}=\prod_{i=n-c+1}^{n}\Bigl(1-\frac ki\Bigr).
\]
\end{proposicion}

\begin{proof}
Fijadas las $n$ muestras, se elige al azar un subconjunto $J$ de $k$ de ellas, uniformemente entre los $\binom nk$ posibles. Hay $\binom{n-c}{k}$ subconjuntos sin ninguna muestra correcta, luego $\binom{n-c}{k}/\binom nk$ es la probabilidad, condicionada a las muestras, de que $J$ no contenga ninguna correcta. Promediando sobre las muestras, la esperanza de este cociente es la probabilidad incondicional de que las $k$ muestras de $J$ sean todas incorrectas. Como $J$ se elige con independencia de los resultados y las muestras son independientes e idénticamente distribuidas, las $k$ muestras de $J$ son $k$ muestras independientes con probabilidad de error $1-p$, y esa probabilidad es $(1-p)^{k}$. Para la segunda igualdad,
\[
\frac{\binom{n-c}{k}}{\binom nk}=\frac{(n-c)!\,(n-k)!}{n!\,(n-c-k)!}=\frac{\prod_{i=n-c+1}^{n}(i-k)}{\prod_{i=n-c+1}^{n}i},
\]
porque $(n-k)!/(n-c-k)!$ es el producto de los $c$ enteros $i-k$ con $i=n-c+1,\dots,n$, y $n!/(n-c)!$ es el producto de los $c$ enteros $i$ del mismo intervalo.
\end{proof}

La forma de producto evita calcular números combinatorios enormes y es la implementación que proponen Chen \emph{et al.}~\cite[§2.1, fig.~3]{chen}, que usan $n=200$ y $k\le 100$.

\begin{ejemplo}[pass@$k$ con diez muestras]\label{ej:passk}
Con $n=10$ y $c=3$: $\widehat{\mathrm{pass}}_1=1-7/10=0{,}3$; $\widehat{\mathrm{pass}}_3=1-\binom73/\binom{10}3=1-35/120=0{,}708$, y $\widehat{\mathrm{pass}}_5=1-\binom75/\binom{10}5=1-21/252=0{,}917$. El estimador ingenuo $1-(1-0{,}3)^{5}$ daría $0{,}832$.
\end{ejemplo}

\section{Calibración}

\begin{definicion}[Calibración, ECE y MCE]\label{def:ece}
Sea un clasificador que devuelve una predicción $\hat Y$ y una confianza $\hat P\in[0,1]$. Está \emph{perfectamente calibrado} si $P(\hat Y=Y\mid\hat P=q)=q$ para todo $q\in[0,1]$. Dadas $n$ predicciones agrupadas en $M$ intervalos de confianza $B_1,\dots,B_M$, sean $\mathrm{acc}_b$ la exactitud y $\mathrm{conf}_b$ la confianza media en $B_b$. El \emph{error de calibración esperado} y el \emph{error de calibración máximo} son
\[
\ECE=\sum_{b=1}^{M}\frac{\abs{B_b}}{n}\,\abs{\mathrm{acc}_b-\mathrm{conf}_b},\qquad \MCE=\max_b\abs{\mathrm{acc}_b-\mathrm{conf}_b}
\]
\cite[§2, ecs.~(1), (3) y~(5)]{guo}.
\end{definicion}

\begin{ejemplo}[Un modelo sobreconfiado]\label{ej:ece}
De 100 respuestas, 50 tienen confianza $0{,}9$ y aciertan el $70\,\%$; 30 tienen confianza $0{,}6$ y aciertan el $60\,\%$, y 20 tienen confianza $0{,}3$ y aciertan el $40\,\%$. Entonces
\[
\ECE=0{,}5\cdot 0{,}2+0{,}3\cdot 0+0{,}2\cdot 0{,}1=0{,}12,
\]
y el MCE es $0{,}2$. La confianza media es $0{,}69$ y la exactitud media, $0{,}61$: el modelo está sobreconfiado.
\end{ejemplo}

El diagrama de fiabilidad representa $\mathrm{acc}_b$ frente a $\mathrm{conf}_b$; un modelo perfectamente calibrado sigue la diagonal. Guo \emph{et al.} observan que las redes profundas modernas están peor calibradas que las antiguas y proponen, entre otros métodos, el \emph{escalado de temperatura}.

\begin{definicion}[Escalado de temperatura]\label{def:temp}
Dados los logits $z$ de un clasificador, el escalado de temperatura devuelve $\softmax(z/T)$, con un único parámetro $T>0$ ajustado minimizando la log-verosimilitud negativa en un conjunto de validación~\cite[§4]{guo}.
\end{definicion}

\begin{proposicion}[El escalado de temperatura no cambia la exactitud]\label{prop:temp}
Para todo $T>0$, $\arg\max_k\softmax(z/T)_k=\arg\max_kz_k$.
\end{proposicion}

\begin{proof}
La exponencial es estrictamente creciente y el denominador de la softmax es común a todas las clases, de modo que $\softmax(z/T)_k>\softmax(z/T)_j$ si y solo si $z_k/T>z_j/T$, que, por ser $T>0$, equivale a $z_k>z_j$.
\end{proof}

El escalado solo cambia la confianza, no la predicción: $T>1$ suaviza las probabilidades y corrige la sobreconfianza. Kadavath \emph{et al.} muestran que los modelos de lenguaje grandes preentrenados están razonablemente bien calibrados en preguntas de opción múltiple con el formato adecuado y que pueden estimar la probabilidad de conocer la respuesta~\cite{kadavath}; el informe técnico de GPT-4 muestra que el posentrenamiento empeoró la calibración de su modelo preentrenado~\cite{gpt4}.

\section{Alucinaciones y reglas de puntuación}

Una \emph{alucinación} es una afirmación falsa formulada con fluidez y seguridad. Kalai \emph{et al.}~\cite{kalai25} sostienen que no es un fallo misterioso, sino una consecuencia estadística del entrenamiento y de la evaluación.
\begin{itemize}[leftmargin=1.4em]
  \item \emph{Preentrenamiento.} Reducen la generación de respuestas válidas a un problema de clasificación binaria («¿es válida esta respuesta?») y prueban que la tasa de error generativo es al menos, aproximadamente, el doble de la tasa de error de esa clasificación~\cite[§3.1, corolario~1]{kalai25}. Para hechos arbitrarios, como la fecha de cumpleaños de una persona poco conocida, la tasa de alucinación de un modelo base está acotada inferiormente por la fracción de hechos que aparecen una sola vez en los datos de entrenamiento~\cite[§3.3.1, teorema~2]{kalai25}, resultado anticipado por Kalai y Vempala~\cite{kalai-vempala}.
  \item \emph{Evaluación.} La mayoría de los conjuntos de evaluación puntúan 1 el acierto y 0 tanto el error como la abstención, de modo que adivinar nunca resta. Un modelo optimizado para esas puntuaciones aprende a no decir nunca «no lo sé».
\end{itemize}

\begin{proposicion}[Umbral de confianza explícito]\label{prop:umbral-t}
Sea $t\in[0,1)$ y una regla de puntuación que da 1 punto por acierto, 0 por abstención y $-t/(1-t)$ por error. Si el modelo acierta con probabilidad $p$, responder tiene mayor puntuación esperada que abstenerse si y solo si $p>t$.
\end{proposicion}

\begin{proof}
La puntuación esperada al responder es $E=p-(1-p)\,t/(1-t)$. Multiplicando por $1-t>0$, $E>0$ si y solo si $p(1-t)-(1-p)t>0$, es decir, $p-pt-t+pt>0$, o sea, $p>t$. La abstención vale~0.
\end{proof}

Kalai \emph{et al.} proponen incluir el umbral en las instrucciones de la evaluación, con valores como $t=0{,}5$, $0{,}75$ o $0{,}9$; la puntuación binaria habitual corresponde a $t=0$~\cite[§4.2]{kalai25}.

\section{\emph{Red teaming} y ataques adversarios}

Una \emph{evaluación de seguridad} mide, antes del despliegue, si un modelo tiene una capacidad peligrosa o si produce conductas dañinas cuando se le pide. El \emph{red teaming} busca activamente entradas que provoquen fallos. Ganguli \emph{et al.} reunieron y analizaron unos $39\,000$ ataques escritos por personas contra modelos de Anthropic~\cite{ganguli}; Perez \emph{et al.} usan un modelo de lenguaje para generar los ataques de forma automática~\cite{perez}.

\begin{definicion}[Ataque por sufijo adversario]\label{def:gcg}
Sea $u$ una petición dañina, $y^{*}=(y^{*}_1,\dots,y^{*}_T)$ el comienzo de una respuesta afirmativa (por ejemplo, «Claro, aquí tienes…») y $s$ un sufijo de tokens. El ataque busca
\[
s^{*}=\arg\min_{s}\ \mathcal{L}(s),\qquad \mathcal{L}(s)=-\sum_{t=1}^{T}\ln p\bigl(y^{*}_t\mid u\oplus s\oplus y^{*}_{<t}\bigr),
\]
donde $\oplus$ denota concatenación~\cite[§2.1]{zou}.
\end{definicion}

El problema es discreto. El algoritmo \emph{Greedy Coordinate Gradient} (GCG) de Zou \emph{et al.} calcula el gradiente de $\mathcal{L}$ respecto de la codificación \emph{one-hot} de cada token del sufijo, toma para cada posición los $k$ tokens con gradiente más negativo como candidatos, evalúa exactamente la pérdida de un lote de sustituciones aleatorias entre esos candidatos y conserva la mejor~\cite[§2.2]{zou}. Optimizado sobre varias peticiones y varios modelos abiertos a la vez, el sufijo resultante se transfiere a modelos cerrados: en su evaluación, la tasa de éxito fue del $86{,}6\,\%$ en GPT-3.5, del $46{,}9\,\%$ en GPT-4, del $47{,}9\,\%$ en Claude~1, del $2{,}1\,\%$ en Claude~2 y del $66{,}0\,\%$ en PaLM-2~\cite[§3.2]{zou}. Son cifras de 2023 sobre versiones ya retiradas; los proveedores corrigen estos ataques una vez publicados.

\begin{figure}[h]
\centering
\begin{tikzpicture}[font=\small,>={Stealth[length=2mm]},
  caja/.style={draw,rounded corners=2pt,minimum height=8mm,align=center,minimum width=20mm}]
  \node[caja] (a) at (0,0) {atacante\\(persona o LLM)};
  \node[caja] (m) at (4.1,0) {modelo\\evaluado};
  \node[caja] (j) at (8.2,0) {juez\\(persona o LLM)};
  \node[caja] (d) at (11.9,0.9) {datos de\\entrenamiento};
  \node[caja] (r) at (11.9,-0.9) {tasa de fallo\\e informe};
  \draw[->] (a) -- node[above] {entradas} (m);
  \draw[->] (m) -- node[above] {respuestas} (j);
  \draw[->] (j) -- (d); \draw[->] (j) -- (r);
  \draw[->] (d.north) -- ++(0,0.4) -| (m.north);
\end{tikzpicture}
\caption{Ciclo de \emph{red teaming}. Los fallos encontrados alimentan el entrenamiento de seguridad y la tasa de fallo informa la decisión de despliegue.}
\label{fig:redteam}
\end{figure}

\section{Estadística de los sucesos raros}

Las conductas peligrosas son, por diseño, raras, y su medida requiere un cuidado estadístico particular.

\begin{proposicion}[Intentos repetidos]\label{prop:intentos}
Si un ataque tiene éxito en cada intento, de forma independiente, con probabilidad $p$, la probabilidad de al menos un éxito en $k$ intentos es $1-(1-p)^{k}$, que es mayor o igual que $1-e^{-kp}$.
\end{proposicion}

\begin{proof}
El suceso complementario es que fallen los $k$ intentos, de probabilidad $(1-p)^{k}$ por la independencia. La desigualdad se sigue de $1-p\le e^{-p}$, consecuencia de la convexidad de la exponencial (su gráfica queda por encima de la recta tangente en~0), elevando ambos miembros, no negativos, a la potencia $k$.
\end{proof}

Con $p=0{,}02$, el ataque triunfa al menos una vez en 100 intentos con probabilidad $1-0{,}98^{100}=0{,}867$. Un ataque que parece raro en una prueba es casi seguro para un atacante persistente.

\begin{proposicion}[Regla del 3]\label{prop:regla3}
Sean $n$ pruebas independientes sin ningún fallo, cada una con probabilidad de fallo $p$ desconocida. La cota superior exacta de confianza $1-\alpha$ para $p$, es decir, el valor $p_u$ con $(1-p_u)^{n}=\alpha$, cumple $p_u\le\ln(1/\alpha)/n$. Para $\alpha=0{,}05$, $\ln 20=2{,}996$ y $p_u\le 3/n$.
\end{proposicion}

\begin{proof}
La probabilidad de no observar ningún fallo es $(1-p)^{n}$, decreciente en $p$; los valores de $p$ incompatibles al nivel $\alpha$ con la observación son los que hacen esa probabilidad menor que $\alpha$, es decir, $p>p_u$. Por la desigualdad $1-p\le e^{-p}$, $\alpha=(1-p_u)^{n}\le e^{-np_u}$; tomando logaritmos, $\ln\alpha\le -np_u$, y por tanto $p_u\le\ln(1/\alpha)/n$.
\end{proof}

Con 300 pruebas sin fallos, la tasa de fallo es menor que $3/300=1\,\%$ con confianza del $95\,\%$; la cota exacta es $1-0{,}05^{1/300}=0{,}99\,\%$. Cero fallos no significa tasa cero.

\section{Marcos de seguridad y conductas ocultas}

Las evaluaciones de seguridad suelen agruparse en cuatro preguntas: si el modelo \emph{puede} hacer algo peligroso (capacidades, como la ayuda en ciberataques o en armas biológicas), si \emph{lo hará} cuando se le pida (salvaguardas, como la tasa de éxito de los ataques), si \emph{persigue} lo que se le pidió (alineamiento: adulación, engaño, sabotaje de tareas) y si \emph{dice lo que representa internamente} (honestidad: calibración y alucinaciones). La \emph{Responsible Scaling Policy} de Anthropic~\cite{rsp} define niveles de seguridad (ASL) que exigen medidas de protección más estrictas cuando las evaluaciones detectan ciertas capacidades; Google DeepMind y OpenAI han publicado marcos análogos.

Hay dos límites conocidos de las evaluaciones de conducta. El primero es que pueden no reflejar el despliegue: contaminación, jueces automáticos sesgados y modelos que distinguen cuándo se les evalúa. El segundo es más profundo: Hubinger \emph{et al.} entrenaron deliberadamente modelos con puertas traseras (por ejemplo, escribir código seguro si el año indicado es 2023 e insertar vulnerabilidades si es 2024) y comprobaron que la conducta persistía tras el ajuste supervisado, el aprendizaje por refuerzo y el entrenamiento adversario, sobre todo en los modelos más grandes; el entrenamiento adversario enseñaba a reconocer mejor el desencadenante y a ocultar la conducta~\cite{hubinger}. Una evaluación de conducta solo ve entradas y salidas; por eso se combina con las herramientas de interpretabilidad de los temas~1 a~5, que miran dentro.

\begin{lecturas}
\item Stanford CS224N~\cite{cs224n}, clase~11 «Benchmarking and Evaluation»; Harvard CS~2881r~\cite{cs2881}, semana~3.
\item Chen \emph{et al.}~\cite{chen}, §2.1 «Functional Correctness».
\item Guo \emph{et al.}~\cite{guo}, §2 y §4; Kadavath \emph{et al.}~\cite{kadavath}.
\item Kalai \emph{et al.}~\cite{kalai25}, §3 y §4.
\item Zou \emph{et al.}~\cite{zou}, §2 y §3; Hubinger \emph{et al.}~\cite{hubinger}.
\item MIT 18.05~\cite{mit1805}, clase~23 «Confidence Intervals for the Mean of Non-normal Data», §3.
\end{lecturas}

\section*{Ejercicios del tema 6}
\addcontentsline{toc}{section}{Ejercicios del tema 6}

\begin{ejercicio}{1}{eje:6-binaria}
Explique por qué un conjunto de evaluación que puntúa 1 el acierto y 0 tanto el error como la abstención incentiva las alucinaciones. Compare la puntuación media de un modelo que adivina con probabilidad de acierto $0{,}1$ con la de uno que se abstiene.
\end{ejercicio}

\begin{ejercicio}{2}{eje:6-ece}
(a) Calcule el ECE de 100 respuestas: 40 con confianza $0{,}95$ y exactitud $0{,}8$; 40 con $0{,}75$ y $0{,}7$; 20 con $0{,}55$ y $0{,}6$. ¿Está el modelo sobreconfiado? (b) Con la regla de la \cref{prop:umbral-t}, ¿le conviene responder a un modelo con $p=0{,}6$ si $t=0{,}5$? ¿Y si $t=0{,}75$? Calcule la puntuación esperada en cada caso.
\end{ejercicio}

\begin{ejercicio}{2}{eje:6-intentos}
Un ataque funciona con probabilidad $0{,}02$ por intento. (a) Calcule la probabilidad de al menos un éxito en 10 y en 100 intentos, y compárela con la aproximación lineal $kp$. (b) En 300 pruebas no hubo ningún éxito: dé la cota superior exacta al $95\,\%$ de la tasa de éxito y la de la regla del~3.
\end{ejercicio}

\begin{ejercicio}{2}{eje:6-passk}
Para un problema se generan $n=20$ soluciones, de las que $c=4$ superan las pruebas. (a) Calcule $\widehat{\mathrm{pass}}_k$ para $k=1$, $5$ y $10$. (b) Calcule $\widehat{\mathrm{pass}}_5$ con la forma de producto de la \cref{prop:passk}. (c) Compare con el estimador ingenuo $1-(1-c/n)^{k}$.
\end{ejercicio}

\begin{ejercicio}{2}{eje:6-mit}
Se mide la exactitud de un modelo en un conjunto de $n$ preguntas con el intervalo conservador de la \cref{prop:ic}. (a) ¿Cuántas preguntas hacen falta para un margen de $0{,}01$ al $95\,\%$? Hágalo primero con $z=2$ y después con $z=1{,}96$. (b) ¿Y al $80\,\%$ ($z_{0{,}1}=1{,}2816$)? (c) Con $n=900$ y una exactitud del $70\,\%$, dé los intervalos al $95\,\%$ y al $80\,\%$. (d) Dos modelos obtienen el $70{,}0\,\%$ y el $72{,}0\,\%$ en $1000$ preguntas distintas cada uno. Estime la desviación típica de la diferencia y decida si es significativa al $95\,\%$.
\fuente{adaptado de MIT 18.05~\cite{mit1805}, primavera de 2022, clase~23, \emph{Board Problem}~1 «Confidence intervals for a binomial» (sondeo electoral sustituido por la exactitud de un modelo; se añade el apartado~(d)).}
\end{ejercicio}

\begin{ejercicio}{3}{eje:6-sesgo}
Sea $c\sim\Bin(n,p)$ y $\hat p=c/n$. (a) Demuestre que $\E\bigl[(1-\hat p)^{2}\bigr]=(1-p)^{2}+p(1-p)/n$ y deduzca que el estimador ingenuo $1-(1-\hat p)^{2}$ subestima pass@2. (b) Con $n=2$, compruebe directamente que el estimador de la \cref{prop:passk} es insesgado para $k=2$. (c) Con $p=0{,}5$ y $n=2$, calcule el valor verdadero de pass@2 y la esperanza del estimador ingenuo.
\end{ejercicio}

\begin{ejercicio}{3}{eje:6-cota-ece}
(a) Demuestre que $\ECE\ge\abs{\overline{\mathrm{conf}}-\overline{\mathrm{acc}}}$, donde $\overline{\mathrm{conf}}$ y $\overline{\mathrm{acc}}$ son la confianza media y la exactitud media de las $n$ predicciones, y dé una condición necesaria y suficiente para la igualdad. (b) Compruebe la desigualdad en el \cref{ej:ece} y en el \cref{eje:6-ece}. (c) Construya un ejemplo con dos intervalos en el que $\overline{\mathrm{conf}}=\overline{\mathrm{acc}}$ y, sin embargo, $\ECE>0$.
\end{ejercicio}
